% !TeX root = ../../Book4.tex
% 纲要标题：### 20.1 射影比较的代数输入与独立入口
% 纲要位置：几何学计划纲要.md，第 4010 行。
% 正文状态：未撰写；本文件目前是框架。

% 纲要内容开始
% 
% 独立入口只重建本章所需的复数域射影版本，不要求先读第五卷全部交换代数、形式函数或模空间理论。
% 
% * 复射影代数对象、概形和结构层的必要构造
% 
%   概形与层对应 V.10.1—10.3、V.13.1—13.3。
% 
% * 代数凝聚层、可逆层、张量与有限呈示
% * 相对仿射构造、Proj 和扭转线丛
% 
%   构造对应 V.13.5、V.14.1—14.2；相对 Spec 不依赖形式函数。
% 
% * 仿射无上同调与有限仿射覆盖的 Čech 模型
% 
%   对应 V.18.1—18.3；不用光滑单位分解替代代数层的无上同调。
% 
% * 射影空间扭转层的上同调及 Serre 生成、消失
% 
%   对应 V.19.1—19.2；解析侧凝聚性取 IV.14.4，高次扭转工具在 IV.20.3 建立。
% 
% * 解析化、局部环映射与态射方向
% 
%   解析化正合性所需的局部环性质在 IV.20.4 就地证明；不得从张量右正合直接推出正合。
% 
% 纲要内容结束

\section{射影比较的代数输入与独立入口}
\label{g4:sec:2001}

待撰写

\subsection{复射影代数对象、概形和结构层的必要构造}
\label{g4:subsec:2001:01}
待撰写

\subsection{代数凝聚层、可逆层、张量与有限呈示}
\label{g4:subsec:2001:02}
待撰写

\subsection{相对仿射构造、Proj 和扭转线丛}
\label{g4:subsec:2001:03}
待撰写

\subsection{仿射无上同调与有限仿射覆盖的 Čech 模型}
\label{g4:subsec:2001:04}
待撰写

\subsection{射影空间扭转层的上同调及 Serre 生成、消失}
\label{g4:subsec:2001:05}
待撰写

\subsection{解析化、局部环映射与态射方向}
\label{g4:subsec:2001:06}
待撰写
% 纲要条目排为编号小节标题，并进入目录与书签；说明段仍保留于纲要注释中。
